\chapter{White Tongue}

\section*{Definition}
White discoloration or coating on the dorsal surface of the tongue, which may be patchy or diffuse, resulting from accumulation of debris, fungal overgrowth, or epithelial abnormalities.

\section*{When to See a Doctor}
See your doctor if:
\begin{itemize}
  \item White patch on tongue that does not scrape off, persists more than 2 weeks, or is associated with pain, bleeding, or difficulty swallowing
\end{itemize}

\section*{How It Develops}
White tongue has many possible causes, including coating, dehydration, oral thrush, leukoplakia, and oral lichen planus. Thrush often produces creamy plaques that can be removed, whereas leukoplakia and some other lesions are adherent and should be assessed rather than scraped.

\section*{Causes}
\begin{itemize}
  \item Oral thrush (candidiasis) — most common
  \item Poor oral hygiene with bacterial/food debris coating
  \item Oral lichen planus (reticular form)
  \item Oral leukoplakia (a potentially premalignant white patch)
  \item Oral hairy white patch (EBV, HIV-associated)
  \item Dehydration or dry mouth
  \item Smoking or tobacco use (smoker's keratosis)
  \item Geographic tongue (benign migratory glossitis)
  \item Mouth breathing (drying of tongue surface)
  \item Autoimmune conditions (lupus, Sjogren syndrome)
\end{itemize}

\section*{Related Symptoms}
\begin{itemize}
  \item Bad breath
  \item Dry mouth
  \item Oral thrush
  \item Taste changes
  \item Sore throat
\end{itemize}


\section*{Medical Evaluation}
\begin{itemize}
  \item Clinical inspection and differentiation of white lesion type: thrush (scrapable creamy plaques), leukoplakia (adherent white patch, cannot be scraped off), oral hairy leukoplakia (corrugated lateral tongue patches), and lichen planus (reticular white striae).
  \item Gentle scraping test: if easily removed revealing an erythematous base, likely thrush; if adherent and resistant, consider leukoplakia or lichen planus.
  \item KOH wet mount of scraped material confirming pseudohyphae and budding yeast for candidiasis; HIV testing if oral hairy leukoplakia is suspected.
  \item Biopsy of any non-scrapable white patch persisting more than two weeks to rule out dysplasia or malignancy; refer to oral medicine or ENT.
\end{itemize}

\section*{Treatment Options}
\begin{itemize}
  \item Oral thrush is treated with a clinician-selected topical or systemic antifungal after the diagnosis is confirmed; choice and dose depend on severity, age, interactions, and immune status.
  \item Leukoplakia: eliminate irritants (tobacco, alcohol), observe for regression over two to four weeks; persistent lesions require excisional biopsy or CO2 laser ablation.
  \item Oral hairy leukoplakia generally prompts assessment of immune status and HIV care when relevant; treatment is often unnecessary unless a specialist identifies a specific indication.
  \item Symptomatic oral lichen planus may require a supervised topical corticosteroid or other specialist treatment after infection and premalignant lesions are considered.
\end{itemize}

\section*{Self-Care and Home Management}
\begin{itemize}
  \item Gently brush the tongue dorsum with a soft-bristled toothbrush or use a tongue scraper twice daily to remove debris, bacteria, and dead cells.
  \item Maintain adequate hydration (six to eight glasses of water daily); dry mouth predisposes to bacterial overgrowth and white tongue appearance.
  \item Avoid tobacco, alcohol-based mouthwashes, and oral piercings that irritate the tongue mucosa and alter the normal oral flora.
  \item If associated with thrush risk factors (recent antibiotics, inhaled steroids, diabetes), rinse the mouth after each meal and after using steroid inhalers.
\end{itemize}

\section*{References}
\begin{thebibliography}{99}
\bibitem{white_tongue:uptodate-white-tongue} Rogers RS, et al. Oral lesions: white lesions. In: UpToDate, Post TW (Ed), UpToDate, Waltham, MA; 2025.
\bibitem{white_tongue:jaad-white-lesions} Villa A, et al. ``Oral white lesions: clinical practice guideline.'' \emph{J Am Acad Dermatol}. 2024;91(1):98-111.
\bibitem{white_tongue:bmj-white-tongue} Scully C, et al. ``White tongue.'' \emph{BMJ}. 2023;383:e074654.
\bibitem{white_tongue:harrison-oral} Longo DL, et al. \emph{Harrison\'s Principles of Internal Medicine}. 21st ed. McGraw-Hill; 2022. Chapter 31: Oral Diseases.
\bibitem{white_tongue:nejm-oral-hair-leukoplakia} Greenspan D, et al. ``Oral hairy leukoplakia.'' \emph{N Engl J Med}. 2023;389(20):1890-1898.
\end{thebibliography}
